%% toc: P2 19 19.8
\section{Proofs for equality with resolution and superposition}
\label{sec:annex-p02-w02b-p02-026}

Terms, clauses, substitutions and the Resolution and Factoring rules are those
of \cref{sec:fol-cardinality-normalforms-04}; the fresh symbol
\(\mathsf{eq}\), the translation \((\cdot)^{\mathsf{eq}}\) and the identity
expansion \(\mathcal N^{=}\) are those of
Convention~\ref{def:p02-w02b-p02-026-reading}.  Three elementary facts are
used throughout.

\emph{(F0) Clause satisfaction.}  Since \(\operatorname{ucl}_{\Lang}(C)\) is
the universal closure of the disjunction of the members of \(C\) (or of
\(\bottom\) when \(C=\varnothing\)), a structure satisfies \(C\) iff every
assignment satisfies some literal \(L\in C\).  No structure satisfies the
empty clause.

\emph{(F1) Coincidence.}  If two assignments agree on the variables of a term
\(t\), then \(t\) has the same value under both; this is an induction on
\(t\).  Hence the truth of a literal \(L\) under an assignment depends only
on the values of the variables of \(L\).

\emph{(F2) Reducts.}  The reduct fact stated at the start of
\cref{sec:annex-p02-w02-p02-028}: a structure and its expansions agree on
formulas of the smaller signature, and every structure of a smaller signature
is the reduct of a structure of a larger one.

\begin{annexitem}{annex:p02-w02b-p02-026-translation}{thm:p02-w02b-p02-026-translation}{Translation lemma}
\begin{proof}
(1)  Substitution acts on the arguments of an atom and leaves its atomic
symbol unchanged, while translation changes only the atomic symbol \(=\) to
\(\mathsf{eq}\) and leaves the arguments unchanged.  Hence
\(((r=t)\theta)^{\mathsf{eq}}=(r\theta\mathrel{\mathsf{eq}}t\theta)
=(r\mathrel{\mathsf{eq}}t)\theta\) and
\((R(\bar t)\theta)^{\mathsf{eq}}=R(\bar t\theta)=R(\bar t)^{\mathsf{eq}}\theta\);
both operations commute with negation, so
\((L\theta)^{\mathsf{eq}}=L^{\mathsf{eq}}\theta\).  Taking pointwise images,
\((C\theta)^{\mathsf{eq}}=\{(L\theta)^{\mathsf{eq}}:L\in C\}
=\{L^{\mathsf{eq}}\theta:L\in C\}=C^{\mathsf{eq}}\theta\).  The variables of a
literal are those of its argument terms, which translation does not change,
so \(\operatorname{Var}(C^{\mathsf{eq}})=\operatorname{Var}(C)\).  For
injectivity: \(L^{\mathsf{eq}}\) has the polarity of \(L\); if its atomic
symbol is \(\mathsf{eq}\), then \(L\) has the same arguments with atomic
symbol \(=\), because \(\mathsf{eq}\notin\Lang\); otherwise \(L\) has the same
atom.  So \(L^{\mathsf{eq}}\) determines \(L\), and consequently
\(C^{\mathsf{eq}}\) determines \(C\) and \(T^{\mathsf{eq}}\) determines
\(T\).

(2)  By (1), \(A^{\mathsf{eq}}\theta=(A\theta)^{\mathsf{eq}}\) and
\(B^{\mathsf{eq}}\theta=(B\theta)^{\mathsf{eq}}\), and by injectivity these are
identical iff \(A\theta=B\theta\).  So the two pairs have the same unifiers.
Being a most general unifier is defined from the set of unifiers and the
instance preorder on substitutions alone, so the two pairs also have the same
most general unifiers.

(3)  By injectivity, \(C\in T\) iff \(C^{\mathsf{eq}}\in T^{\mathsf{eq}}\).
For a variable renaming \(\pi\), \((C\pi)^{\mathsf{eq}}=C^{\mathsf{eq}}\pi\)
by (1), and two clauses are renamed apart iff their translations are, since
the variable sets agree.  Let \(C_k\) follow from earlier \(C_i,C_j\) by
Resolution: \(E=C_i\pi\) and \(F=C_j\rho\) are renamed apart, \(A\in E\),
\(\neg B\in F\), \(\theta\) is a most general unifier of \(A,B\), and
\(C_k=\bigl((E\setminus\{A\})\cup(F\setminus\{\neg B\})\bigr)\theta\).  Then
\(E^{\mathsf{eq}}=C_i^{\mathsf{eq}}\pi\) and
\(F^{\mathsf{eq}}=C_j^{\mathsf{eq}}\rho\) are renamed apart,
\(A^{\mathsf{eq}}\in E^{\mathsf{eq}}\), \(\neg B^{\mathsf{eq}}\in F^{\mathsf{eq}}\),
\(\theta\) is a most general unifier of \(A^{\mathsf{eq}},B^{\mathsf{eq}}\) by
(2), and, since injectivity gives
\((E\setminus\{A\})^{\mathsf{eq}}=E^{\mathsf{eq}}\setminus\{A^{\mathsf{eq}}\}\)
and likewise for \(F\), (1) yields
\[
 C_k^{\mathsf{eq}}
 =\bigl((E^{\mathsf{eq}}\setminus\{A^{\mathsf{eq}}\})\cup
   (F^{\mathsf{eq}}\setminus\{\neg B^{\mathsf{eq}}\})\bigr)\theta ,
\]
the Resolution conclusion from \(C_i^{\mathsf{eq}},C_j^{\mathsf{eq}}\).  If
\(C_k=C_i\theta\) by Factoring on distinct literals \(L_1,L_2\in C_i\) of the
same polarity with atoms \(A,B\) and most general unifier \(\theta\), then
\(L_1^{\mathsf{eq}},L_2^{\mathsf{eq}}\) are distinct members of
\(C_i^{\mathsf{eq}}\) of that polarity, \(\theta\) is a most general unifier
of their atoms by (2), and \(C_k^{\mathsf{eq}}=C_i^{\mathsf{eq}}\theta\).
Conversely, let \(C_k^{\mathsf{eq}}\) follow by Resolution from
\(C_i^{\mathsf{eq}}\pi\) and \(C_j^{\mathsf{eq}}\rho\) on a positive
\(A'\in C_i^{\mathsf{eq}}\pi=(C_i\pi)^{\mathsf{eq}}\) and a negative
\(\neg B'\in(C_j\rho)^{\mathsf{eq}}\) with most general unifier \(\theta\).
Then \(A'=A^{\mathsf{eq}}\) for a positive \(A\in C_i\pi\) and
\(\neg B'=\neg B^{\mathsf{eq}}\) for a negative \(\neg B\in C_j\rho\);
\(C_i\pi,C_j\rho\) are renamed apart, \(\theta\) is a most general unifier of
\(A,B\) by (2), and the Resolution conclusion \(R\) from \(C_i\pi,C_j\rho\) on
\(A,\neg B\) with \(\theta\) satisfies \(R^{\mathsf{eq}}=C_k^{\mathsf{eq}}\) by
the computation above, so \(R=C_k\) by injectivity.  Factoring is reversed in
the same way.  Applying these facts to every member of the sequence proves
(3).

(4)  \(\mathcal N\) and \(\mathcal N^{=}\) have the same domain and interpret
every symbol of \(\Lang\) alike, so every term has the same value in both under
\(s\).  Then \(\mathcal N,s\entails r=t\) iff
\(\sem r^{\mathcal N}_s=\sem t^{\mathcal N}_s\) iff the pair of values lies in
the identity relation, which interprets \(\mathsf{eq}\) in \(\mathcal N^{=}\),
iff \(\mathcal N^{=},s\entails r\mathrel{\mathsf{eq}}t\).  For a relation atom
the two structures agree, and negation is preserved on both sides.  The claim
for clauses follows by (F0), since \(\mathcal N\) and \(\mathcal N^{=}\) have
the same assignments.
\end{proof}
\end{annexitem}

\begin{annexitem}{annex:p02-w02b-p02-026-quotient}{thm:p02-w02b-p02-026-quotient}{The equality axioms capture identity}
\begin{proof}
\emph{Left to right.}  Let \(\mathcal N\) satisfy \(T\).  By
Lemma~\ref{thm:p02-w02b-p02-026-translation}(4), \(\mathcal N^{=}\) satisfies
\(T^{\mathsf{eq}}\).  Every clause of \(\mathrm{Eq}(\Sigma)\) is satisfied by
\(\mathcal N\), by (F0) and the following checks under an arbitrary assignment
\(s\).  (E1): \(s(x)=s(x)\).  (E2) and (E3): if the negative literals are all
false, then \(s(x)=s(y)\) (and \(s(y)=s(z)\)), so \(s(y)=s(x)\) (and
\(s(x)=s(z)\)).  (E4): if \(s(x_i)=s(y_i)\) for all \(i\), the argument tuples
of \(f^{\mathcal N}\) coincide, so the two values coincide.  (E5): if
\(s(x_i)=s(y_i)\) for all \(i\) and \(R^{\mathcal N}\) holds of
\((s(x_1),\ldots,s(x_n))\), it holds of the same tuple
\((s(y_1),\ldots,s(y_n))\).  By
Lemma~\ref{thm:p02-w02b-p02-026-translation}(4) again, \(\mathcal N^{=}\)
satisfies \(\mathrm{Eq}(\Sigma)^{\mathsf{eq}}\).

\emph{Right to left.}  Let \(\mathcal N'\) be an
\((\Lang\cup\{\mathsf{eq}\})\)-structure with domain \(D'\) satisfying
\(T^{\mathsf{eq}}\cup\mathrm{Eq}(\Sigma)^{\mathsf{eq}}\), and write \(\sim\) for
the relation interpreting \(\mathsf{eq}\).

\emph{Step 1: \(\sim\) is a congruence for \(\Sigma\).}  The variables
displayed in each clause of \(\mathrm{Eq}(\Sigma)\) are pairwise distinct, so
for any elements of \(D'\) there is an assignment sending those variables to
them.  Applying (F0) to the translated clauses under such assignments gives,
for all elements of \(D'\): \(a\sim a\) by (E1); \(a\sim b\) implies \(b\sim a\)
by (E2); \(a\sim b\) and \(b\sim c\) imply \(a\sim c\) by (E3); for an
\(n\)-ary \(f\in\Sigma\) with \(n\geq1\), \(a_i\sim b_i\) for all \(i\) implies
\(f^{\mathcal N'}(\bar a)\sim f^{\mathcal N'}(\bar b)\) by (E4); and for an
\(n\)-ary \(R\in\Sigma\) with \(n\geq1\), \(a_i\sim b_i\) for all \(i\) and
\(R^{\mathcal N'}(\bar a)\) imply \(R^{\mathcal N'}(\bar b)\) by (E5).  So
\(\sim\) is an equivalence relation; write \([a]\) for the class of \(a\).

\emph{Step 2: the quotient.}  Fix \(d_0\in D'\).  Let \(\mathcal Q\) be the
\(\Lang\)-structure with domain \(\{[a]:a\in D'\}\), which is nonempty, and
\begin{itemize}
  \item \(c^{\mathcal Q}=[c^{\mathcal N'}]\) for a constant \(c\in\Sigma\),
        and \(c^{\mathcal Q}=[d_0]\) for any other constant;
  \item \(f^{\mathcal Q}([a_1],\ldots,[a_n])=[f^{\mathcal N'}(a_1,\ldots,a_n)]\)
        for \(f\in\Sigma\) of arity \(n\geq1\), and the constant value
        \([d_0]\) for any other function symbol;
  \item \(R^{\mathcal Q}([a_1],\ldots,[a_n])\) iff
        \(R^{\mathcal N'}(a_1,\ldots,a_n)\) for \(R\in\Sigma\) (for \(n=0\),
        \(R^{\mathcal Q}\) is the truth value \(R^{\mathcal N'}\)), and
        \(R^{\mathcal Q}\) empty (false for \(n=0\)) for any other relation
        symbol.
\end{itemize}
These do not depend on the representatives: for \(f\) by the congruence
property of Step~1, and for \(R\) because \([a_i]=[b_i]\) for all \(i\) gives
\(R^{\mathcal N'}(\bar a)\Rightarrow R^{\mathcal N'}(\bar b)\) by Step~1 and
the reverse implication by Step~1 applied with symmetry.  As every
\(\Lang\)-structure, \(\mathcal Q\) interprets \(=\) as identity.

\emph{Step 3: terms.}  Call a term, literal or clause a \(\Sigma\)-term,
\(\Sigma\)-literal or \(\Sigma\)-clause when every constant, function and
relation symbol in it belongs to \(\Sigma\).  For an assignment \(s\) in
\(\mathcal N'\) let \([s]\) be the assignment \(v\mapsto[s(v)]\) in
\(\mathcal Q\).  For every \(\Sigma\)-term \(t\),
\(\sem t^{\mathcal Q}_{[s]}=\bigl[\sem t^{\mathcal N'}_s\bigr]\): for a
variable this is the definition of \([s]\); for a constant of \(\Sigma\) it is
the definition of \(c^{\mathcal Q}\); and for \(f(t_1,\ldots,t_n)\) with
\(f\in\Sigma\) the induction hypothesis and the definition of
\(f^{\mathcal Q}\) give
\(f^{\mathcal Q}([\sem{t_1}^{\mathcal N'}_s],\ldots,[\sem{t_n}^{\mathcal N'}_s])
=[\sem{f(t_1,\ldots,t_n)}^{\mathcal N'}_s]\).

\emph{Step 4: literals.}  For every \(\Sigma\)-literal \(L\),
\(\mathcal Q,[s]\entails L\) iff \(\mathcal N',s\entails L^{\mathsf{eq}}\).
For \(r=t\): by Step~3, \(\sem r^{\mathcal Q}_{[s]}=\sem t^{\mathcal Q}_{[s]}\)
iff \([\sem r^{\mathcal N'}_s]=[\sem t^{\mathcal N'}_s]\) iff
\(\sem r^{\mathcal N'}_s\sim\sem t^{\mathcal N'}_s\) iff
\(\mathcal N',s\entails r\mathrel{\mathsf{eq}}t\).  For \(R(t_1,\ldots,t_n)\)
with \(R\in\Sigma\): by Step~3 and the definition of \(R^{\mathcal Q}\), it
holds in \(\mathcal Q\) under \([s]\) iff \(R^{\mathcal N'}\) holds of the
values \(\sem{t_i}^{\mathcal N'}_s\).  Negation is preserved on both sides.

\emph{Step 5: \(\mathcal Q\) satisfies \(T\).}  Let \(C\in T\); it is a
\(\Sigma\)-clause by the hypothesis on \(\Sigma\).  Let \(\beta\) be an
assignment in \(\mathcal Q\), and let \(z_1,\ldots,z_k\) be the finitely many
variables of \(C\).  Choose \(a_i\in\beta(z_i)\) for \(i=1,\ldots,k\); these
are finitely many choices from nonempty sets, made one at a time, so no choice
principle is needed.  Let \(s(z_i)=a_i\) and \(s(v)=d_0\) for every other
variable.  Since \(\mathcal N'\) satisfies \(C^{\mathsf{eq}}\), (F0) gives
\(L\in C\) with \(\mathcal N',s\entails L^{\mathsf{eq}}\), so
\(\mathcal Q,[s]\entails L\) by Step~4.  As \([s](z_i)=[a_i]=\beta(z_i)\) and
the variables of \(L\) are among the \(z_i\), (F1) gives
\(\mathcal Q,\beta\entails L\).  So every assignment satisfies a literal of
\(C\), and \(\mathcal Q\) satisfies \(C\) by (F0).  Hence \(T\) is
satisfiable.
\end{proof}
\end{annexitem}

\begin{annexitem}{annex:p02-w02b-p02-026-eq-resolution}{thm:p02-w02b-p02-026-eq-resolution}{Soundness and refutational completeness of resolution with equality axioms}
\begin{proof}
Put \(T=(S\cup\mathrm{Eq}_S)^{\mathsf{eq}}
=S^{\mathsf{eq}}\cup\mathrm{Eq}_S^{\mathsf{eq}}\).  It is a finite
equality-free set of clauses of \(\Lang\cup\{\mathsf{eq}\}\), a signature that
contains a constant of \(\Lang\) (Convention~\ref{def:p02-w02b-p02-026-reading},
item~2), so Convention~\ref{def:p02-w02-p02-028-standing} holds for it.

(1)\(\Leftrightarrow\)(2).  Lemma~\ref{thm:p02-w02b-p02-026-quotient} with
\(T:=S\) and \(\Sigma:=\Sigma_S\).

(2)\(\Leftrightarrow\)(3).  By Theorem~\ref{thm:p02-w02-p02-028-clause-form}
applied to \(T\), \(T\) is unsatisfiable iff some finite set \(G'\) of ground
instances of members of \(T\) is propositionally unsatisfiable, with ground
terms built from the symbols occurring in \(T\) and one constant when these
include none.  As \(\mathsf{eq}\) is a relation symbol, those terms use only
symbols of \(\Sigma_S\) and possibly one constant of \(\Lang\).  By
Lemma~\ref{thm:p02-w02b-p02-026-translation}(1), a ground instance of
\(C^{\mathsf{eq}}\) is \(C^{\mathsf{eq}}\gamma=(C\gamma)^{\mathsf{eq}}\), so
the finite sets of ground instances of members of \(T\) are exactly the sets
\(G^{\mathsf{eq}}\) with \(G\) a finite set of ground instances of members of
\(S\cup\mathrm{Eq}_S\).  The map \(A\mapsto A^{\mathsf{eq}}\) is injective on
ground atoms.  If an atom valuation \(w\) satisfies \(G^{\mathsf{eq}}\), then
\(v(A)\defeq w(A^{\mathsf{eq}})\) gives every ground literal \(L\) the value
that \(w\) gives \(L^{\mathsf{eq}}\), so \(v\) satisfies \(G\).  Conversely, if
\(v\) satisfies \(G\), let \(w(A^{\mathsf{eq}})\defeq v(A)\), which is well
defined by injectivity, and \(w\) arbitrary on other atoms; then \(w\)
satisfies \(G^{\mathsf{eq}}\).  So \(G\) is propositionally unsatisfiable iff
\(G^{\mathsf{eq}}\) is.

(2)\(\Leftrightarrow\)(4).  By Theorem~\ref{thm:p02-w02-p02-028-resolution},
(1)\(\Leftrightarrow\)(3), applied to \(T\), \(T\) is unsatisfiable iff
Resolution with Factoring derives the empty clause from \(T\).  Every clause
in a derivation from \(T\) is free of \(=\): the members of \(T\) are, and a
Resolution or Factoring conclusion consists of substitution instances of
literals of its parents, with the same atomic symbols.  A clause free of
\(=\) is \(K^{\mathsf{eq}}\) for the \(\Lang\)-clause \(K\) obtained by
writing \(=\) for \(\mathsf{eq}\).  Hence every derivation from \(T\) is
\(C_1^{\mathsf{eq}},\ldots,C_n^{\mathsf{eq}}\) for a sequence of
\(\Lang\)-clauses, which by Lemma~\ref{thm:p02-w02b-p02-026-translation}(3)
is a derivation from \(S\cup\mathrm{Eq}_S\); and the translation of every
derivation from \(S\cup\mathrm{Eq}_S\) is a derivation from \(T\).  Since
\(\varnothing^{\mathsf{eq}}=\varnothing\) and translation is injective, one of
them contains the empty clause iff the other does.

The soundness half (4)\(\Rightarrow\)(1) can also be read off directly: by
the soundness theorem of \cref{sec:fol-cardinality-normalforms-04}, a
derivation of the empty clause from \(S\cup\mathrm{Eq}_S\) shows that no
\(\Lang\)-structure satisfies \(S\cup\mathrm{Eq}_S\), and every
\(\Lang\)-structure satisfies \(\mathrm{Eq}_S\) (the left-to-right part of the
proof of Lemma~\ref{thm:p02-w02b-p02-026-quotient}), so no
\(\Lang\)-structure satisfies \(S\).
\end{proof}
\end{annexitem}

\begin{annexitem}{annex:p02-w02b-p02-026-bare}{thm:p02-w02b-p02-026-bare}{What resolution without the axioms refutes}
\begin{proof}
\(S^{\mathsf{eq}}\) is a finite equality-free set of clauses of
\(\Lang\cup\{\mathsf{eq}\}\), which contains a constant.  By the argument of
the last part of the preceding proof, with \(S\) in place of
\(S\cup\mathrm{Eq}_S\), Resolution with Factoring derives the empty clause
from \(S\) iff it does so from \(S^{\mathsf{eq}}\), and by
Theorem~\ref{thm:p02-w02-p02-028-resolution} that happens iff
\(S^{\mathsf{eq}}\) is unsatisfiable.  If \(\mathcal N\) satisfied \(S\), then
\(\mathcal N^{=}\) would satisfy \(S^{\mathsf{eq}}\) by
Lemma~\ref{thm:p02-w02b-p02-026-translation}(4); so unsatisfiability of
\(S^{\mathsf{eq}}\) implies that of \(S\).  The failure of the converse is
Corollary~\ref{thm:p02-w02b-p02-026-failures}(1).
\end{proof}
\end{annexitem}

\begin{annexitem}{annex:p02-w02b-p02-026-failures}{thm:p02-w02b-p02-026-failures}{Two incomplete treatments of equality}
\begin{proof}
In both items the clauses are ground, so by (F0) a structure satisfies such a
clause iff some member is true in it.  Symbols of the signature that do not
occur in the clauses are interpreted arbitrarily, which is possible by (F2)
and does not affect the clauses.

(1)  \emph{\(S_0\) is unsatisfiable.}  If \(\mathcal N\) satisfies
\(\{a=b\}\), then \(a^{\mathcal N}=b^{\mathcal N}\), so
\(f^{\mathcal N}(a^{\mathcal N})=f^{\mathcal N}(b^{\mathcal N})\); then
\(P(f(a))\) and \(P(f(b))\) have the same truth value in \(\mathcal N\), and
\(\{P(f(a))\}\) and \(\{\neg P(f(b))\}\) cannot both hold.  \emph{No
refutation.}  By Theorem~\ref{thm:p02-w02b-p02-026-bare} it suffices to show
that \(S_0^{\mathsf{eq}}=\bigl\{\{a\mathrel{\mathsf{eq}}b\},\{P(f(a))\},
\{\neg P(f(b))\}\bigr\}\) is satisfiable.  Let \(\mathcal M\) have domain
\(\{0,1,2,3\}\), \(a^{\mathcal M}=0\), \(b^{\mathcal M}=1\),
\(f^{\mathcal M}(0)=2\), \(f^{\mathcal M}(1)=3\),
\(f^{\mathcal M}(2)=f^{\mathcal M}(3)=2\), \(P^{\mathcal M}=\{2\}\) and
\(\mathsf{eq}^{\mathcal M}=\{(0,1)\}\).  Then \(a\mathrel{\mathsf{eq}}b\) holds
since \((0,1)\in\mathsf{eq}^{\mathcal M}\); \(P(f(a))\) holds since
\(f^{\mathcal M}(0)=2\in P^{\mathcal M}\); and \(\neg P(f(b))\) holds since
\(f^{\mathcal M}(1)=3\notin P^{\mathcal M}\).

(2)  \emph{\(S_1\) is unsatisfiable.}  If \(a^{\mathcal N}=b^{\mathcal N}\),
then \(P(a)\) and \(P(b)\) have the same truth value, so \(\{P(a)\}\) and
\(\{\neg P(b)\}\) cannot both hold.  \emph{No refutation.}  \(\Sigma'\)
contains no function symbol of positive arity and no relation symbol, so
\(\mathrm{Eq}(\Sigma')\) consists of (E1), (E2) and (E3) only.  By
Theorem~\ref{thm:p02-w02b-p02-026-bare} applied to the finite set
\(S_1\cup\mathrm{Eq}(\Sigma')\), it suffices that
\((S_1\cup\mathrm{Eq}(\Sigma'))^{\mathsf{eq}}\) is satisfiable.  Let
\(\mathcal M'\) have domain \(\{0,1\}\), \(a^{\mathcal M'}=0\),
\(b^{\mathcal M'}=1\), \(P^{\mathcal M'}=\{0\}\) and
\(\mathsf{eq}^{\mathcal M'}=\{0,1\}\times\{0,1\}\).  The translated (E1),
(E2) and (E3) hold under every assignment, because their positive literal
\(x\mathrel{\mathsf{eq}}x\), \(y\mathrel{\mathsf{eq}}x\) or
\(x\mathrel{\mathsf{eq}}z\) is true for every pair; \(a\mathrel{\mathsf{eq}}b\)
holds; \(P(a)\) holds as \(0\in P^{\mathcal M'}\); and \(\neg P(b)\) holds as
\(1\notin P^{\mathcal M'}\).
\end{proof}
\end{annexitem}

\begin{annexitem}{annex:p02-w02b-p02-026-superposition-sound}{thm:p02-w02b-p02-026-superposition-sound}{Soundness of the superposition rules}
\begin{proof}
\emph{Positions.}  The positions of a term are finite sequences of positive
integers: the empty sequence \(\epsilon\) is a position of every term, with
\(s|_\epsilon=s\) and \(s[u]_\epsilon=u\); if \(s=f(s_1,\ldots,s_n)\) and
\(p\) is a position of \(s_i\), then \(i.p\) is a position of \(s\), with
\(s|_{i.p}=s_i|_p\) and
\(s[u]_{i.p}=f(s_1,\ldots,s_{i-1},s_i[u]_p,s_{i+1},\ldots,s_n)\).
Two facts follow by induction on \(p\).  (P1) If \(p\) is a position of
\(s\), it is a position of \(s\sigma\), \((s\sigma)|_p=(s|_p)\sigma\) and
\((s[u]_p)\sigma=(s\sigma)[u\sigma]_p\): the case \(\epsilon\) is immediate,
and the case \(i.p\) applies the induction hypothesis to \(s_i\), since
\(f(s_1,\ldots,s_n)\sigma=f(s_1\sigma,\ldots,s_n\sigma)\).  (P2) If \(u\) and
\(u'\) have the same value in \(\mathcal N\) under an assignment \(\beta\),
then so do \(v[u]_p\) and \(v[u']_p\) for every position \(p\) of \(v\): the
case \(\epsilon\) is the hypothesis, and in the case \(i.p\) the two terms
apply the same function \(f^{\mathcal N}\) to argument values that agree
everywhere, by the induction hypothesis at place \(i\).

On the sheet, a clause \(C\lor L\) is the clause \(C\cup\{L\}\), and clauses
are taken modulo variable renaming, which preserves satisfaction (the
soundness theorem of \cref{sec:fol-cardinality-normalforms-04} records both
directions).  Let \(\mathcal N\) satisfy the premises of an inference and let
\(\sigma\) be its substitution.  By Lemma~\ref{lem:clause-substitution},
applied to each premise separately, \(\mathcal N\) satisfies every premise
instantiated by \(\sigma\); this needs no renaming apart.  Fix an assignment
\(\beta\).  By (F0) it suffices to find a literal of the conclusion true
under \(\beta\), and if \(\beta\) satisfies a literal of \(C\sigma\) or
\(D\sigma\), that literal belongs to the conclusion.  Assume therefore that it
satisfies none.

\emph{Superposition.}  The premises are \(C\cup\{l=r\}\) and
\(D\cup\{s[l']_p\diamond t\}\), where \(p\) is the position of the
distinguished occurrence \(l'\) in \(s\) and \(\diamond\) is \(=\) or its
negation; the conclusion is \((C\cup D\cup\{s[r]_p\diamond t\})\sigma\), and
\(l\sigma=l'\sigma\).  By the assumption, \(\beta\) satisfies
\(l\sigma=r\sigma\) and \((s[l']_p\diamond t)\sigma\).  By (P1),
\((s[l']_p)\sigma=(s\sigma)[l'\sigma]_p=(s\sigma)[l\sigma]_p\) and
\((s[r]_p)\sigma=(s\sigma)[r\sigma]_p\).  Since \(l\sigma\) and \(r\sigma\)
have the same value under \(\beta\), (P2) shows that \((s[l']_p)\sigma\) and
\((s[r]_p)\sigma\) do.  The truth of a literal \(u\diamond t\sigma\) under
\(\beta\) depends only on the values of \(u\) and \(t\sigma\), so \(\beta\)
satisfies \((s[r]_p\diamond t)\sigma\).  If the distinguished occurrence lies
in the right-hand side of the literal instead, the same argument applies with
the two sides exchanged.

\emph{Equality Resolution.}  The premise is \(C\cup\{\neg(s=t)\}\), the
conclusion \(C\sigma\), and \(s\sigma=t\sigma\).  The literal
\(\neg(s\sigma=t\sigma)\) is false under \(\beta\), because identical terms
have the same value; so \(\beta\) satisfies a literal of \(C\sigma\), contrary
to the assumption.  Hence this case cannot arise, and the conclusion holds.

\emph{Equality Factoring.}  The premise is \(C\cup\{s=t,\ u=v\}\), the
conclusion \((C\cup\{\neg(t=v),\ u=v\})\sigma\), and \(s\sigma=u\sigma\).  If
\(\beta\) satisfies \(u\sigma=v\sigma\) or \(\neg(t\sigma=v\sigma)\), it
satisfies the conclusion.  Otherwise \(t\sigma\) and \(v\sigma\) have the same
value and \(u\sigma=v\sigma\) is false; since the premise instance holds,
\(s\sigma=t\sigma\) is true.  Then the value of \(u\sigma\), which is the
value of \(s\sigma\), equals the value of \(t\sigma\) and hence of
\(v\sigma\), so \(u\sigma=v\sigma\) is true, a contradiction.

In every case \(\beta\) satisfies the conclusion; as \(\beta\) was arbitrary,
\(\mathcal N\) satisfies it.  The ordering, eligibility and selection
conditions were not used.  If the rules derive the empty clause from \(S\),
induction along the derivation shows that every model of \(S\) satisfies
every derived clause, including the empty clause, which no structure
satisfies by (F0); so \(S\) has no model.
\end{proof}
\end{annexitem}
